% TOP_PROBLEM: {"id": 3140, "title": "Sidorenko's Conjecture", "category_id": 3, "category": {"id": 3, "name": "graph_theory", "display_name": "Graph Theory", "description": "Problems involving graphs, networks, and their properties.", "slug": "graph-theory", "order_index": 3, "created_at": "2026-07-31T15:26:25.671Z"}, "status": "open", "problem_number": "OPG-60039", "set_id": 12}
% FIRST_POSTED: 2026-10-01
% ENTRY_KIND: statement_only
% UnsolvedMath ID 3140; source catalogue code OPG-60039
% !TeX program = lualatex
% Definition and sources only; this file is not an LLM solution attempt.
\documentclass[11pt,a4paper]{article}
\usepackage[margin=25mm]{geometry}
\usepackage{fontspec}
\setmainfont{lmroman10-regular.otf}[BoldFont=lmroman10-bold.otf,
 ItalicFont=lmroman10-italic.otf,BoldItalicFont=lmroman10-bolditalic.otf]
\usepackage{amsmath,amssymb,mathtools}
\usepackage{unicode-math}
\setmathfont{latinmodern-math.otf}
\AtBeginDocument{\let\setminus\smallsetminus}
\usepackage{xurl,hyperref}
\hypersetup{colorlinks=true,urlcolor=blue}
\setcounter{secnumdepth}{0}
\setlength{\parindent}{0pt}
\setlength{\parskip}{5pt}
\setlength{\emergencystretch}{3em}
\begin{document}
\section*{Sidorenko's Conjecture}\label{3140}
{\small UnsolvedMath ID 3140 · OPG-60039 · Graph Theory · OpenGarden}
\subsection{Definitions and mathematical statement}
Let $H$ and $G$ be finite simple undirected graphs, with
$N=|V(G)|\geq1$. A graph homomorphism from $H$ to $G$ is a function
$f:V(H)\to V(G)$ for which $f(u)f(v)\in E(G)$ whenever
$uv\in E(H)$. The function need not be injective; distinct
nonadjacent vertices of $H$ may have the same image. Let
$\operatorname{hom}(H,G)$ count these functions, and set
\[
 t(H,G)=\frac{\operatorname{hom}(H,G)}{N^{|V(H)|}},
 \qquad p(G)=\frac{2|E(G)|}{N^2}=t(K_2,G).
\]
Thus $t(H,G)$ is the probability that a uniformly chosen function
on the vertices preserves all edges.

A graph is bipartite if its vertices can be partitioned into two
sets with no edges inside either set. Sidorenko's conjecture states
that, for every bipartite $H$ and every nonempty $G$,
\[
 t(H,G)\geq p(G)^{|E(H)|}.
\]
Equivalently, multiply the right side by $N^{|V(H)|}$ to obtain
the lower bound on $\operatorname{hom}(H,G)$. Isolated vertices of
$H$ are allowed and contribute freely chosen images. If $H$ has
no edges the right side is taken to be $1$, including when $p(G)=0$.
The normalization uses $N^2$, because images of the two ends of
an edge are chosen independently and in order. It is not the usual
density with denominator $\binom N2$.
\subsection{Short English statement}
Does every bipartite pattern have at least as many edge-preserving vertex maps as its normalized edge density predicts?
\subsection{Sources}
\begin{itemize}
\item[S1] ULAM.AI and the UnsolvedMath Contributors, supplied problems.json, record 3140 (OPG-60039), OpenGarden collection. Catalogue identity and source transcription; CC BY 4.0.
\url{https://huggingface.co/datasets/ulamai/UnsolvedMath}.
\item[S2] Open Problem Garden, Sidorenko's Conjecture. Original problem and discussion; the mathematical formulation below is expanded from the supplied source record.
\url{http://www.openproblemgarden.org/op/sidorenkos_conjecture}.
\end{itemize}
\end{document}
